% Integration fragment: no documentclass.
% Uses the main manuscript's notation and Lemma lem:semiprime55.
% Do not interpret bounded multiplicity as arbitrary arithmetic weights.

\subsection{The same upper sieve for finitely many prime factors}
\label{sec:primeproducts55}

The upper-weight factorization established in
Lemma~\ref{lem:semiprime55} is a property of the modulus coefficients.
It can also be used for the switched product sequences below.
We record the additional argument, particularly the coprime main
mass and the multiplicity of the factorizations.

Put $a=100/1149$ and $K_0=\lfloor1/a\rfloor=11$.
An admissible product multiset here means a finite positive sum
of labelled prime-tuple families
\[
 m=p_1\cdots p_k,\qquad 1\le k\le K_0,
 \qquad p_j\ge N^a,\qquad m\le N,
\]
whose domains are specified by finitely many ordering and monomial
inequalities in the $p_j$ and $N$. Labels distinguish different
specified-prefix choices for the same product.
Their coefficients are fixed bounded nonnegative numbers.
In particular, the resulting multiplicity $w(m)$ is bounded
independently of $N$. This does not mean that an arbitrary
bounded arithmetic function $w(m)$ is permitted: the tuple-domain
description is needed for separation of variables.

\begin{lemma}[Derived upper sieve for finite prime-product multisets]
\label{lem:primeproducts55}
Let $\mathcal M_N$ be an admissible product multiset as above,
and write
\[
 X_N=\sum_m w(m).
\]
Assume its mass has a limit
$X_N=(J+o(1))N/\log N$, evaluated with its stated labels and
multiplicities. Then
\begin{equation}\label{eq:primeproduct-prime-upper}
 \sum_{\substack{m<N\\N-m\ {\rm prime}}}w(m)
 \le\left(\frac{80}{11}J+o(1)\right)\HN .
\end{equation}
More generally the same conclusion applies to a finite positive
sum of such domains and to positive smooth outer approximations,
with their own total mass $X_N$.

This lemma applies to all five switched multisets defined by
$Y_7,Y_8,Y_9,Y_{10},Y_{11}^{\rm full}$ in
Section~\ref{sec:repaired-chen-weight}. If $U_i=8J_i$ denotes
the corresponding upper constant computed there with coefficient
eight, then
\begin{equation}\label{eq:switch-55}
 \Upsilon_i\le\left(\frac{10}{11}U_i+o(1)\right)\HN
 \quad(7\le i\le10),\qquad
 \Upsilon_{11}^{\rm full}
       \le\left(\frac{10}{11}U_{11}^{\rm full}+o(1)\right)\HN.
\end{equation}
Here the last mass uses the complete range
$N^h\le p_4<N^d$.
\end{lemma}

\begin{proof}
Fix $0<\varepsilon_0<1/100$ and put
$D=N^{11/20-\varepsilon_0}$ and $z=N^{1/5}$.
Use precisely the Iwaniec modified upper weight $\rho_D^+$
and its normalized box-atom expansion constructed in
Lemma~\ref{lem:semiprime55}. In particular
\[
 |\rho_D^+(d)|\le C_\eta,\qquad d\mid P(z),
\]
where $\eta$ is the fixed construction precision.
Its factorization works for all medium lengths
$U\in[x^{2\zeta},x^{1/3+\zeta/2}]$.
The positive margins from the two cases
$t_1\ge21/200$ and $t_1<21/200$ depend only on modulus
factorization, not on the sequence being counted.
Use
\[
 \eta\ll\epsilon\ll\zeta\ll\varepsilon_0,\qquad
 12\zeta<a/2,\qquad \kappa\ll\epsilon .
\]
All parameters and $K_0$ are fixed before $N\to\infty$.

\paragraph{Small products, repeated primes, and smooth domains.}
Products below $N^{1-\kappa}$ cost
$O(N^{1-\kappa})$ in the prime count because $w$ is bounded.
On $m\le CN$, the products having a repeated prime at least
$N^a/2$ cost at most
\[
 C' N\sum_{\ell\ge N^a/2,\ \ell\ {\rm prime}}\ell^{-2}
       \ll N^{1-a}
\]
for bounded multiplicity. Thus actual repeated primes can be
removed before using strictly ordered prime tuples, with an
additive error of this size in the desired upper bound.
This removal does not identify distinct labelled prefix choices.

Choose a single nonnegative smooth outer approximation of each
tuple domain in the exponent variables
$\log p_j/\log N$ and the product variable $m/N$.
Preserve a separate lower-cutoff factor which is one when
$p_j\ge N^a$ and zero when $p_j\le N^a/2$.
Thus smoothing never lowers the individual prime-factor bound
to $N^{a-\delta}$; all products in the outer family still have
every prime factor at least $N^a/2$.
This constant-factor boundary enlargement has vanishing width
in exponent coordinates.

A smooth dyadic partition of unity in each coordinate gives
$O((\log N)^k)$ rectangles, with bounded overlap and uniformly
bounded normalized derivatives, and sums back to the original
global majorant. For each fixed smoothing precision $\delta>0$,
the prime number theorem and partial summation evaluate its total
mass. On the compact exponent domains $t_j\ge a+o(1)$, the
prime-product integral has bounded density; the finite monomial
and ordering boundaries have measure zero. Hence the excess
mass is $o_\delta(1)N/\log N$, where $o_\delta(1)\to0$
as $\delta\downarrow0$. This is one global mass estimate;
it is not repeated for each dyadic rectangle.
The same statement follows directly from the rough-number mass
evaluation in the five applications below.
A component supported only on a boundary can instead be included
in a shrinking boundary strip.

\paragraph{Distribution with the exact coprime mass.}
For a smooth $k$-tuple weight $W$ on one rectangle, put
\[
 X_W^{(d)}
   =\sum_{\substack{p_1,\ldots,p_k\ {\rm prime}\\
                    (p_1\cdots p_k,d)=1}}
                         W(p_1,\ldots,p_k).
\]
We first prove, on scales $N^{1-\kappa}/O(1)\le x\ll N$,
\begin{equation}\label{eq:primeproduct-coprime-error}
 \sum_{\substack{d\mid P(z)\\(d,N)=1}}\rho_D^+(d)
 \left\{
  \sum_{\substack{p_1,\ldots,p_k\ {\rm prime}\\
            p_1\cdots p_k\equiv N\pmod d}}W(p_1,\ldots,p_k)
       -\frac{X_W^{(d)}}{\varphi(d)}
 \right\}
 \ll_A N(\log N)^{-A}.
\end{equation}

Apply the $k_{\rm HB}=3$ Heath--Brown identity to each of the
$k$ prime slots, at that slot's own scale, initially with
$\Lambda$ weights. There are at most $6k\le66$ variables,
of which at most $3k$ carry a $\mu$ coefficient.
Retain the condition
$P^-(p_1\cdots p_k)>x^{1/(\log\log x)^3}$ before expansion.
Consequently every nonunit variable is rough at that cutoff.
Each $\mu$ variable has length at most the cube root of
its own original factor scale, and in particular at most
$x^{1/3+o(1)}$.
Before Mellin separation, discard every dyadic block that cannot
intersect its original prime-slot interval. If the original slot has
scale $R_i$, the product of that group's at most six dyadic lower
endpoints is at least $R_i/2^6\ge N^a/O(1)$. This group-scale bound
persists after separating the smooth weights, including at points
outside the original curved domain.

If an individual variable has length
\[
 x^{2\zeta}\le U\le x^{1/3+\zeta/2},
\]
take that individual variable as the Type II factor and group
all the others. Its coefficient is $1$ or $\mu$, with a possible
logarithm. Lemma~\ref{lem:rough-sw}, with $\xi=2\zeta$,
verifies its complete rough Siegel--Walfisz input.
The other coefficient is divisor-bounded, with fixed order
depending on $K_0$. The normalized three-factor modulus
coefficients satisfy $L_{23}$, Proposition~5.2, by the
factorization already proved in Lemma~\ref{lem:semiprime55}.
The truncation of divisor-bounded coefficients and their
logarithmic losses are handled exactly as in that lemma.

Suppose no individual variable is in this middle interval.
All $\mu$ variables are then shorter than $x^{2\zeta}$.
Every original prime group has at most six variables and
original product at least $N^a/2$. It must therefore contain
a variable longer than $x^{1/3+\zeta/2}$: otherwise its
product would be at most $x^{12\zeta+o(1)}<N^a/2$.
On the other hand there are at most two such long variables
globally, because the entire product is comparable with $x$.
It follows that this case is impossible for $k\ge3$.

For $k=2$, there is exactly one long variable in each original
group. Both have coefficient $1$ or a logarithm, and the
remaining ten variables have product at most $x^{20\zeta+o(1)}$.
Apply $L_{23}$, Proposition~3.2:
\[
 N_1,N_2>x^{1/3+\zeta/2}>x^{3\epsilon},
 \qquad
 Q_{\rm modulus}\le x^{3/5}
       <(x/x^{20\zeta+o(1)})^{2/3-3\epsilon}.
\]
For $k=1$, either there is one long variable and the other
five have product at most $x^{10\zeta+o(1)}$, so Lemma~3.4
applies, or there are two long variables and the remaining
four have product at most $x^{8\zeta+o(1)}$, so
Proposition~3.2 applies. These include all alternatives.

The separate original product cutoffs are retained using
the $k$-variable Mellin transform of $W$, truncated at
polylogarithmic frequencies. A chosen individual Type II
variable acquires only a twist $n^{i\tau}$; partial summation
preserves its Siegel--Walfisz estimates with a polylogarithmic
loss. The two long variables, when present, are handled by
two partial summations, using the arbitrary-subinterval form
explicit in the proof of Maynard's Proposition~5.2,
\href{https://arxiv.org/pdf/2006.07088v1}
{arXiv:2006.07088v1}, which is the input for $L_{23}$,
Proposition~3.2.
The one-long-variable case allows the same partial summation
in Lemma~3.4.
No Siegel--Walfisz property of an arbitrary convolution or
arbitrary-coefficient extension of the two-long-variable
estimate is assumed.

For every local scale use the fixed residue parameter
$\alpha=1+3\kappa$. Then $N<x^\alpha<x^{1+\epsilon}$
for sufficiently small $\kappa$ and sufficiently large $N$.
Changes of exponent coordinates are $O(\kappa)$ and are
absorbed by the fixed factorization margins.
Rough modulus atoms with at most two large-prime slots have
support below $N^{2\nu/3+O(\eta)}$, $\nu=11/20-\varepsilon_0$.
They are covered by ordinary Bombieri--Vinogradov for $k=1$.
For $k\ge2$, select one original prime-interval factor and
group the other $k-1$; both have lengths between fixed
positive powers of the product scale. The generalized
Bombieri--Vinogradov input \eqref{eq:geh}, with its verified
Siegel--Walfisz factor, applies. Smooth Mellin twists cost
only logarithmic powers.

As in Lemma~\ref{lem:semiprime55}, Mellin tails are controlled
after summing over divisors of $N-m$, using
\[
 \sum_{\substack{m\ll N\\m\ne N}}
      \tau_C(m)\tau_C(|N-m|)
        \ll N(\log N)^{C'},
\]
which follows by Cauchy--Schwarz and divisor second moments.
Any isolated $m=N$ term is $O(DN^{o(1)})$.
Proper prime powers introduced by the $\Lambda$ identities
have total mass at most $N^{1-a/2+o(1)}$; after the divisor
argument their cumulative progression error still has
this power saving. Their mean error uses
$\sum_d|\rho_D^+(d)|/\varphi(d)\ll_\eta(\log N)^{C_\eta}$.
Remove logarithmic prime weights by partial summation.
This proves \eqref{eq:primeproduct-coprime-error}.
Increasing $A$ absorbs the number of dyadic rectangles,
Mellin losses, and the fixed number of prime-factor patterns.

\paragraph{Replacing the coprime mass: a power-saving bound.}
Here $X_W^{(d)}$ cannot be identified with $X_W$:
some prime factors may be below $z$.
The exact correction is nevertheless negligible even after
the modulus weights have been summed in absolute value.
For a prime $\ell$,
\[
 \sum_{\substack{d\mid P(z)\\\ell\mid d}}
     \frac{|\rho_D^+(d)|}{\varphi(d)}
 \le \frac{C_\eta}{\ell-1}
          \sum_{e\le D/\ell}\frac1{\varphi(e)}
 \ll_\eta \frac{\log N}{\ell-1}.
\]
Indeed $d$ is squarefree, so $d=\ell e$ with $(e,\ell)=1$
and $\varphi(d)=(\ell-1)\varphi(e)$; dropping the residual
support restrictions only increases this positive bound.
After summing the smooth tuples, let $w_{\rm out}(m)$ denote
their nonnegative bounded total multiplicity and $X_{\rm out}$
their total mass. The union bound therefore gives
\begin{align*}
 &\sum_{\substack{d\mid P(z)\\(d,N)=1}}
    \frac{|\rho_D^+(d)|}{\varphi(d)}
                      (X_{\rm out}-X_{\rm out}^{(d)})\\
 &\qquad\ll_\eta
    (\log N)\sum_m w_{\rm out}(m)
          \sum_{\ell\mid m}\frac1{\ell-1}\\
 &\qquad\ll_{\eta,K_0}
    X_{\rm out}N^{-a}\log N
       \ll N^{1-a}(\log N)^{C}.
\end{align*}
Every such $\ell$ is at least $N^a/2$ and each product has at
most $K_0$ distinct factors. This proves the claimed uniform
power saving. It is valid for positive bounded labelled
multiplicities and does not lose a factor $D$.
If the original switched set excludes products sharing a prime
factor with $N$, the raw mass of those products is separately bounded by
\[
 O\left(N\sum_{\substack{\ell\mid N\\\ell\ge N^a/2}}\frac1\ell\right)
       =N^{1-a+o(1)}.
\]
This deletion is therefore also harmless; it is distinct from the
coprime-main-mass correction just estimated.

We may now replace $X_W^{(d)}$ by the total mass throughout
\eqref{eq:primeproduct-coprime-error}.
Apply the upper linear sieve with density
$g_N(\ell)=1/(\ell-1)$ for $\ell\nmid N$ and zero otherwise.
The uniform small-prime presieve and the full product $V_N(z)$
are those established in Lemma~\ref{lem:semiprime55}.
For $s=5\nu\in(1,3)$ one has $F(s)=2e^\gamma/s$, and hence
\[
 S\le\left(\frac4\nu+o(1)\right)
                    \frac{\CN X_{\rm out}}{\log N}.
\]
If the outer approximation reaches products just above $N$,
sieve $|N-m|$; every desired representation is still counted.
The cases $0<N-m\le z$ in the desired prime count contribute
$O(N^{1/5})$ for bounded
multiplicity. First let $N\to\infty$, then remove the smooth
outer approximation and the sieve precision within the fixed
margins, and finally let $\varepsilon_0\downarrow0$.
This proves \eqref{eq:primeproduct-prime-upper}.

\paragraph{The five switched multisets and the rough quotient.}
In each of $Y_7,Y_8,Y_9,Y_{10},Y_{11}^{\rm full}$, keep the
specified primes with exactly their original labels and ordering.
Factor the rough quotient uniquely as
\[
 u=q_1\cdots q_j,\qquad q_1<\cdots<q_j,
\]
after the repeated-prime exception above is removed.
Split into the finitely many possible values of $j$.
This is a canonical ordered factorization of the quotient;
there is no additional division by $j!$ and no deletion of
different specified-prefix choices for the same full product.
All quotient primes are at least the prescribed roughness
threshold, and in these five domains that threshold is at least
$N^a$. Thus the entire product has at most $K_0$ factors.
The resulting domains satisfy the finite monomial-inequality
description used above, including the moving quotient threshold
$(N/(p_1p_2))^{1/2}$ in $Y_9$.

The case $u=1$ has negligible mass. The five possible specified
prefix products are bounded respectively by
\[
 N^{2/3},\quad N^{2/3},\quad N^{(1+d)/2},
                  \quad N^{4h},\quad N^{3h+d},
\]
all smaller than $N$ by a fixed power.
Likewise, specified factors dividing $N$ or any repeated
prime affect the original prime count only by a power-saving
error. No restriction of the fourth-prime range is introduced:
for $Y_{11}^{\rm full}$ it remains
$N^a\le p_1<p_2<p_3<N^h\le p_4<N^d$.

Reassembling the canonical quotient factorizations gives exactly
the original rough-number mass, with its prefix multiplicities.
Its Buchstab/prime-number-theorem evaluation is therefore
unchanged. In the notation of the main calculation,
\[
 X_i=(J_i+o(1))\frac{N}{\log N},
 \qquad J_i=\frac{U_i}{8}.
\]
The treatment of the Buchstab argument-one boundary by a thin
strip is preserved; smoothing does not remove its mass silently.
Applying the just-proved upper sieve to each of these masses
gives \eqref{eq:switch-55}.
\end{proof}

\paragraph{Scope.}
This is a derived extension of the modulus-atom and smooth
product argument in Lemma~\ref{lem:semiprime55}.
It retains the full domains and tuple multiplicities of the
five switching terms. It changes only their analytic upper
coefficient from $8$ to $80/11$.
It makes no assertion of distribution for arbitrary bounded
arithmetic reweightings or of bibliographic priority.
